\documentclass{beamer}
\usetheme{Boadilla}


\title{Orchestration Technique of Intelligent Microservices on the Edge Computing \\in \\ High-Resource-Constraint Devices}
%btitle{Using beamer for presentations}
\author{Caio Jos\'e Borba Vilar Guimar\~aes}
\institute{UFCG}

\date{\today}

\begin{document}    

\begin{frame}
\titlepage
\end{frame}

\begin{frame}
    \frametitle{Outline}
    \tableofcontents
\end{frame}

\section{Objectives}
\subsection{Main Objective}
\begin{frame}
\frametitle{Main Objective}
To develop, evaluate, and validate a technique for orchestrating intelligent microservices within a distributed and collaborative Edge Computing (EC) infrastructure composed of resource-constrained devices.
    
\end{frame}

\begin{frame}
    \begin{list}{$\bullet$}{}
        \item To design a clear topology architecture for the distributed and collaborative Edge Computing infrastructure;
        \item To implement a collaborative and distributed scheduling system that intelligently instantiate microservices on connected devices;
        \item To create methods and algorithms for intelligent microservices orchestration;
        \item To test and evaluate these solutions both in simulated environments and real IoT setups, covering scenarios typical of Smart Environments and IoT applications.
    \end{list}

\end{frame}

\subsection{Specific Objectives}
\begin{frame}
\frametitle{Specific Objectives}
\begin{list}{$\bullet$}
    \item Study paradigms of orchestration techniques and protocols for collaborative distribution among contrained devices;
    \item Design and develop a collaborative microservices orchestration technique tailored for high-restriction devices in distributed and cooperative environments, particularly in smart environments;
    \item Simulate existing state-of-the-art techniques in smart environments to evaluate their performance in terms of transmission rate, energy consumption, and scalability in microservices deployment;
    \item Develop an infrastructure with high-restriction devices and diverse communication technologies (wired and wireless) to test selected techniques in scenarios closely resembling real-world IoT environments;
\end{list}
\end{frame}

\begin{frame}
\frametitle{Specific Objectives}
\begin{list}{$\bullet$}
    \item Implement and evaluate AI microservices (Machine Learning and Deep Learning models)—assessing accuracy, inference latency, and provisioning time in both simulated and real IoT environments;
    \item Develop new methods and algorithms for orchestration and cooperative control of intelligent microservices, optimizing the use of limited resources and mitigating the limitations of existing techniques;
    \item Analyze and compare the performance of simulation and real-world implementations, refining techniques and protocols based on empirical data;
    \item Ensure the developed protocols are feasible and conformant, applying formal analysis methods to verify their viability considering system constraints like memory and processing power;
    \item Optimize existing algorithms and protocols by improving efficiency and reducing resource consumption, adapting based on feedback from initial implementations and tests 
\end{list}
\end{frame}


\end{document}
